%-------------------------
% Resume in LaTeX -- French master, mirrors templates/cv_base.tex block for block
% (same entries, bullets and dates; tests/test_cv_base_fr.py checks it).
% Based on Jake Gutierrez template, customized for Hongming Fang
% ONE-PAGE MASTER — condensed from cv_base_fr.tex (2 bullets per experience, 1 per project); same facts, fewer details
%------------------------

\documentclass[a4paper,12pt]{article}

\usepackage{latexsym}
\usepackage{titlesec}
\usepackage{marvosym}
\usepackage[usenames,dvipsnames]{color}
\usepackage{verbatim}
\usepackage{enumitem}
\usepackage[hidelinks]{hyperref}
\usepackage{fancyhdr}
\usepackage[french]{babel}
\frenchbsetup{StandardLists=true}
\usepackage{tabularx}
\usepackage{fontawesome5}
\usepackage{multicol}
\usepackage{ragged2e}
\usepackage{needspace}
\usepackage[
  top=0.3in,
  bottom=0.5in,
  left=0.55in,
  right=0.55in
]{geometry}

\setlength{\multicolsep}{-3.0pt}
\setlength{\columnsep}{-1pt}
\input{glyphtounicode}

%----------HEADER / FOOTER----------
\pagestyle{fancy}
\fancyhf{} % clear all header and footer fields
\fancyfoot{}
\renewcommand{\headrulewidth}{0pt}
\renewcommand{\footrulewidth}{0pt}

\urlstyle{same}

% A page that ends early (the next block cannot start there) is stretched, through the glue
% before sections and entries, instead of leaving a blank; the last page is never stretched.
\flushbottom
\vbadness=10000
\clubpenalty=10000 \widowpenalty=10000   % never leave a single line of a bullet alone on a page
\justifying
\setlength{\tabcolsep}{0in}

%----------SECTION FORMATTING----------
\titleformat{\section}{
  \vspace{-4pt}\scshape\raggedright\large\bfseries
}{}{0em}{}[\color{black}\titlerule \vspace{-5pt}]
\titlespacing*{\section}{0pt}{3.5ex plus 30pt minus .2ex}{2.3ex plus .2ex}
% A section heading never ends a page alone: it moves with the first lines of its content.
% Projects need more (their heading and description cannot split), see \needspace before them.
\let\oldsection\section
\renewcommand{\section}[1]{\needspace{5\baselineskip}\oldsection{#1}}

% Ensure that generated pdf is machine readable/ATS parsable
\pdfgentounicode=1

%-------------------------
% Custom commands
\newcommand{\resumeItem}[1]{
  \item\small{
    {#1\vspace{-1pt}}
  }
}

\newcommand{\classesList}[4]{
    \item\small{
        {#1 #2 #3 #4 \vspace{-2pt}}
  }
}

% An entry heading stays with its first bullet (at least two lines, see \clubpenalty) without needing a fixed gap
\newcommand{\resumeSubheading}[4]{
  \vspace{-2pt plus 12pt}\item
    \begin{tabular*}{1.0\textwidth}[t]{l@{\extracolsep{\fill}}r}
      \textbf{#1} & \textbf{#2} \\
      {#3} & \textnormal{#4} \\
    \end{tabular*}\vspace{-7pt}\nopagebreak[4]
}

\newcommand{\resumeSubSubheading}[2]{
    \item
    \begin{tabular*}{0.97\textwidth}{l@{\extracolsep{\fill}}r}
      \textit{\small#1} & \textit{\small #2} \\
    \end{tabular*}\vspace{-7pt}
}

\newcommand{\resumeProjectHeading}[2]{
    \item
    \begin{tabular*}{1.001\textwidth}{l@{\extracolsep{\fill}}r}
      \small#1 & \textbf{\small #2}\\
    \end{tabular*}\vspace{-7pt}
}

\newcommand{\resumeProjectHeadingFourItemResearch}[4]{
    \needspace{6\baselineskip}\item
    \begin{tabular*}{1.001\textwidth}{l@{\extracolsep{\fill}}r}
      #1 & \textbf{#2}\\
      {#3} & \textnormal{#4} \\
    \end{tabular*}\vspace{-7pt}
}

\newcommand{\resumeSubItem}[1]{\resumeItem{#1}\vspace{-4pt}}

\renewcommand\labelitemi{$\vcenter{\hbox{\tiny$\bullet$}}$}
\renewcommand\labelitemii{$\vcenter{\hbox{\tiny$\bullet$}}$}

\newcommand{\resumeSubHeadingListStart}{\begin{itemize}[leftmargin=0.0in, label={}]}
\newcommand{\resumeSubHeadingListEnd}{\end{itemize}}
\newcommand{\resumeItemListStart}{\begin{itemize}}
\newcommand{\resumeItemListEnd}{\end{itemize}\vspace{-5pt}}

\newcommand{\projectdesc}[1]{%
  {\fontsize{11pt}{12pt}\selectfont #1}\vspace{-6pt}%
}

%-------------------------------------------
%%%%%%  RESUME STARTS HERE  %%%%%%%%%%%%%%%%%%%%%%%%%%%%


\begin{document}

%----------HEADING----------
\vspace*{-0.15in}
\begin{center}
    {\LARGE \scshape Hongming Fang} \\[2pt]
    \small Paris, France $|$ (+33) 07 75 77 55 06 $|$
    \href{mailto:hongming.marius.fang@gmail.com}{hongming.marius.fang@gmail.com} $|$
    \href{https://linkedin.com/in/hongming-fang}{\raisebox{-0.2\height}\faLinkedin\ linkedin.com/in/hongming-fang} \\
    {Ingénieur Machine Learning (CDI/CDD), disponible immédiatement — Île-de-France, ouvert à la mobilité}
    \vspace{-5pt}
\end{center}

%SUMMARY-BEGIN
\noindent\small{Ingénieur Machine Learning avec une expérience concrète de l'ingénierie LLM et de la segmentation d'images médicales, acquise lors de stages ML à Paris et à Pékin.}\par
\vspace{-4pt}
%SUMMARY-END

%-----------EDUCATION-----------
\section{FORMATION}
  \resumeSubHeadingListStart
    \resumeSubheading
    {Master Informatique, parcours DIGIT, spécialité Image}{09/2024 -- 09/2026}
    {Sorbonne Université}{Paris, France}
    \resumeItemListStart
     \resumeItem{\textit{Enseignements principaux} : Traitement d'images, Algorithmes d'informatique graphique 3D, Science des données, Architecture des ordinateurs, Imagerie biologique et médicale, Reconnaissance des formes pour l'analyse et l'interprétation des images, Techniques avancées de vision par ordinateur}
     \resumeItem{Diplômé avec mention Très Bien (16,3/20 de moyenne générale).}
    \resumeItemListEnd

    \resumeSubheading
      {Diplôme d'Ingénieur en Électronique et Informatique}{09/2017 -- 06/2024}
      {École Centrale de Pékin, Université de Beihang}{Pékin, Chine}
    \resumeItemListStart
        \resumeItem{Diplôme d'ingénieur français accrédité par la CTI, issu d'un programme conjoint entre l'Université de Beihang et le Groupe des Écoles Centrales, avec licence et master chinois en sciences.}
    \resumeItemListEnd
  \resumeSubHeadingListEnd

%-----------PROFESSIONAL EXPERIENCE-----------
\section{EXPÉRIENCE PROFESSIONNELLE}
  \resumeSubHeadingListStart

    \resumeSubheading
      {DiliTrust}{03/2026 -- 08/2026}
      {Stagiaire ingénieur Machine Learning}{Paris, France}
      \resumeItemListStart
        \resumeItem{Conception d'un cadre de prompts structuré enrichi de métadonnées pour l'extraction d'entités juridiques multilingues avec Qwen3.5-4B (+1,6 point d'extraction precision, +8,4 points de context accuracy), et d'une boucle automatisée où un LLM-as-a-judge affine les prompts. Résultat : une context accuracy de 96,8\,\% et un temps d'inférence réduit de 40,8\,\% par rapport au modèle de production.}
        \resumeItem{Fine-tuning de Qwen3.5-4B par SFT et apprentissage par renforcement (GRPO) : l'accuracy in-distribution dépasse de 8 points celle du modèle de base, avec la meilleure généralisation aux labels jamais vus de toutes les variantes fine-tunées.}
        \resumeItem{Conception et implémentation d'un pipeline de génération de métadonnées qui exploite l'historique d'annotations des utilisateurs pour améliorer automatiquement les consignes d'extraction.}
      \resumeItemListEnd

    \resumeSubheading
      {DeepWise}{03/2023 -- 09/2023}
      {Stagiaire en algorithmes d'IA (imagerie médicale)}{Pékin, Chine}
      \resumeItemListStart
        \resumeItem{Amélioration du modèle de segmentation de la table d'examen et des objets externes grâce à l'ajout de scans CTA de l'aorte nouvellement annotés, avec des scores Dice de 0,88/0,72. Après le diagnostic d'une baisse de performance sur une anatomie out-of-distribution, entraînement d'un nouveau modèle avec des scans CTA de la tête supplémentaires, utilisé pour pré-annoter plus de 1\,000 CT sans injection (moins de 5\,\% de corrections manuelles), puis réentraînement d'un modèle unifié sur des données CTA et CT, avec des scores Dice de 0,93/0,76 et une bonne généralisation à des régions du corps jamais vues. Ajout d'un post-traitement qui supprime les petites composantes connexes isolées.}
      \resumeItemListEnd

    \resumeSubheading
      {Orange Labs International Center Beijing}{06/2022 -- 09/2022}
      {Stagiaire en algorithmes d'IA (modélisation de la parole auto-supervisée)}{Pékin, Chine}
      \resumeItemListStart
        \resumeItem{Reproduction du pipeline complet de génération de parole sans texte (GSLM) à partir des articles de recherche, avec des encoders auto-supervisés et des vocoders neuronaux, incluant l'entraînement d'un encoder HuBERT en distributed data parallel (DDP) sur un corpus bilingue de 30 Go construit sur mesure. Analyse des modes d'échec en anglais et en français qui limitaient la généralisation au français.}
      \resumeItemListEnd

  \resumeSubHeadingListEnd

%-----------PROJECTS & RESEARCH-----------
\needspace{8\baselineskip}
\section{PROJETS ET RECHERCHE}
    \vspace{-5pt}
    \resumeSubHeadingListStart

    \resumeProjectHeadingFourItemResearch
    {\textbf{Job Hunter : recherche d'emploi et adaptation de CV automatisées par LLM}}{08/2026 -- 09/2026}
    {Projet personnel}{Paris, France}
    \resumeItemListStart
      \resumeItem{Mise en place de la collecte de plus de 3\,000 offres auprès de 17 sources, où une source en échec n'interrompt pas les autres, avec déduplication inter-sources (comparaison exacte et vérification par LLM) et pré-filtrage par un score à base de règles. Ajout d'un LLM-as-a-judge qui note chaque offre par rapport au profil du candidat et déclenche l'adaptation automatique du CV : le LLM se limite à sélectionner parmi les expériences et les bullet points existants du CV, et ne peut donc rien inventer. Automatisation par des tâches planifiées et un watchdog qui relance les exécutions manquées, avec un tableau de bord FastAPI sur SQLite pour le suivi.}
    \resumeItemListEnd

    \resumeProjectHeadingFourItemResearch
    {\textbf{Analyse d'images et de nuages de points pour l'histoire de l'art}}{11/2025 -- 02/2026}{Sorbonne Université, projet PRAT}{Paris, France}
    \resumeItemListStart
      \resumeItem{Construction d'un pipeline de recalage de nuages de points coarse-to-fine sous Open3D (sous-échantillonnage par voxels, alignement global RANSAC, raffinement ICP) pour aligner des scans photogrammétriques et des scans laser d'un relief en bois de 12 mètres. Comparaison d'un descripteur géométrique traditionnel (FPFH) et d'un descripteur appris (SpinNet) : SpinNet, pré-entraîné uniquement sur des scènes d'intérieur, réussit sur 5 fragments sur 7 où FPFH échoue, pour 6,6 fois plus de temps d'exécution.}
    \resumeItemListEnd

    \resumeProjectHeadingFourItemResearch
    {\textbf{Reproduction de NeRF sous PyTorch}}{04/2025 -- 05/2025}
    {Sorbonne Université, projet du cours IG3D}{Paris, France}
    \resumeItemListStart
      \resumeItem{Reproduction de NeRF (reconstruction de scènes 3D à partir d'images) sous PyTorch, validée par des ablation studies sur des jeux de données standard et appliquée à un jeu de données Blender personnalisé, où l'ajustement des sampling bounds a corrigé des distorsions géométriques.}
    \resumeItemListEnd

    \resumeProjectHeadingFourItemResearch
    {\textbf{What Makes a Great City to Live in France}}{02/2025 -- 05/2025}
    {Sorbonne Université, projet du cours DALAS}{Paris, France}
    \resumeItemListStart
      \resumeItem{Collecte de données de plusieurs sources (web scraping et API) sur 20 grandes villes françaises, puis conversion en indicateurs par habitant comparables. Visualisation de la distribution et des corrélations des variables, puis combinaison d'une ACP, d'un hierarchical clustering et d'une analyse par SHAP, montrant que la qualité de vie urbaine résulte d'un équilibre entre activité économique et facteurs liés au mode de vie.}
    \resumeItemListEnd

    \resumeProjectHeadingFourItemResearch
    {\textbf{Apprentissage profond pour la classification des polypes colorectaux}}{02/2025 -- 05/2025}
    {Sorbonne Université, projet DIGIT}{Paris, France}
    \resumeItemListStart
      \resumeItem{Classification des polypes colorectaux (adénomateux ou hyperplasiques) par CNN VGG-19 sur le jeu de données clinique POLAR, après prétraitement et restructuration du jeu pour garantir la cohérence des labels. Application du transfer learning avec data augmentation pour atténuer le déséquilibre des classes, ce qui a fait passer le F1 des polypes hyperplasiques de 0,148 à 0,351.}
    \resumeItemListEnd

    \resumeProjectHeadingFourItemResearch
     {\textbf{Interpretability Based Dialog Policy Optimization of TOD}}{12/2022 -- 06/2024}{Université de Beihang}{Mémoire de Master}
    \resumeItemListStart
        \resumeItem{Modélisation d'un task-oriented dialogue à l'aide de modèles causaux, puis entraînement de la dialogue policy par apprentissage par renforcement hors ligne (BCQ) sur MultiWOZ 2.0. Résultats compétitifs (Inform 96,3\,\%, Success 83,4\,\%), avec une robustesse et une interprétabilité supérieures à celles des modèles de référence.}
    \resumeItemListEnd

    \resumeSubHeadingListEnd

%-----------SKILLS-----------
\section{COMPÉTENCES}
 \begin{itemize}[leftmargin=0.2in, label={}]
    \small{\item{
    \textbf{Compétences techniques :} Python, C/C++ ; PyTorch, NumPy, pandas, scikit-learn, OpenCV, Open3D ; FastAPI, SQLite, API REST, pytest ; Linux, Git, CI/CD, Docker\\
    \textbf{IA générative et systèmes agentiques :} prompt engineering, RAG, fine-tuning de LLM (LoRA, SFT, GRPO), pipelines agentiques LangGraph, structured LLM output, vLLM\\
    \textbf{Évaluation et monitoring des LLM :} évaluation LLM-as-a-judge, conception de failure taxonomies, diagnostic des hallucinations et de la généralisation, optimisation du coût et de la latence d'inférence\\
    \textbf{Deep Learning :} CNN, Transformers, self-supervised learning, transfer learning, distributed training (DDP)\\
    \textbf{Vision par ordinateur :} traitement d'images, segmentation, recalage de nuages de points 2D/3D, classification\\
    \textbf{Langues :} mandarin (langue maternelle), anglais (courant), français (niveau professionnel)
    }}
 \end{itemize}

\end{document}
